% =============================================================================
% Kapitel 1: Einleitung
% -----------------------------------------------------------------------------
% Die Einleitung fuehrt in das Thema ein und endet mit Forschungsfrage,
% Methodik und Aufbau. Sie ist das einzige Kapitel OHNE Kapitel-Einleitungssatz
% vor der ersten \subsection (alle anderen Kapitel brauchen 2--4 einleitende
% Saetze zwischen \section und erster \subsection -- siehe FOM-Leitfaden).
% =============================================================================
\section{Einleitung}
\label{sec:einleitung}

% --- 1.1 Problemstellung und Relevanz ---
\subsection{Problemstellung und Relevanz}
\label{subsec:problemstellung}

Hier beschreibst du das Problem und warum es relevant ist. Beginne mit einem
konkreten Aufhänger und arbeite dich zur wissenschaftlichen Lücke vor. Jede
Tatsachenbehauptung wird belegt. Ein indirektes (sinngemäßes) Zitat sieht so
aus:\autocite[Vgl.][S.~15]{mustermann2024}
Ein direktes Zitat steht in Anführungszeichen und ohne \enquote{Vgl.}:
\enquote{Wörtlich übernommener Satz.}\autocite[][S.~16]{mustermann2024}

Fachbegriffe werden bei der ersten Verwendung eingeführt. Danach genügt die
Abkürzung, z.\,B. \ac{KMU}. Bei der ersten Nennung schreibt \verb|\ac| den
Begriff automatisch aus.

% --- 1.2 Forschungsfrage ---
\subsection{Forschungsfrage}
\label{subsec:forschungsfrage}

Aus der Problemstellung ergibt sich die folgende Forschungsfrage:

\begin{quote}
  \itshape
  Hier steht deine eine, präzise Forschungsfrage. Sie ist spezifisch,
  beantwortbar und im Umfang der Arbeit zu bewältigen.
\end{quote}

% --- 1.3 Methodik und Aufbau ---
\subsection{Methodik und Aufbau der Arbeit}
\label{subsec:methodik}

Hier erläuterst du, mit welcher Methode du die Frage beantwortest
(z.\,B. Literaturanalyse, empirische Erhebung, Fallstudie) und wie die Arbeit
aufgebaut ist. Verweise auf die folgenden Kapitel mit \verb|\ref|:
Kapitel~\ref{sec:grundlagen} legt die Grundlagen, Kapitel~\ref{sec:hauptteil}
analysiert den Kern der Fragestellung, und Kapitel~\ref{sec:fazit} fasst die
Ergebnisse zusammen.

Am Ende dieses Abschnitts nennst du den verwendeten FOM-Leitfaden, an dem sich
die formale Gestaltung orientiert (FOM-Regel).
